% Generated draft: requires mathematical and source review.
\documentclass[11pt]{article}
\usepackage[margin=25mm]{geometry}
\usepackage{fontspec}
\setmainfont{Latin Modern Roman}
\usepackage{amsmath,amssymb,mathtools,unicode-math}
\setmathfont{Latin Modern Math}
\usepackage{enumitem,needspace,xurl,hyperref}
\setlength{\parindent}{0pt}
\setlength{\parskip}{5pt}
\setlength{\emergencystretch}{3em}
\newcommand{\N}{\mathbb{N}}
\newcommand{\Z}{\mathbb{Z}}
\newcommand{\Q}{\mathbb{Q}}
\newcommand{\R}{\mathbb{R}}
\newcommand{\C}{\mathbb{C}}
\newcommand{\F}{\mathbb{F}}
\newcommand{\A}{\mathbb{A}}
\newcommand{\PP}{\mathbb{P}}
\newcommand{\E}{\mathbb{E}}
\newcommand{\eps}{\varepsilon}
\newcommand{\dd}{\,\mathrm{d}}
\begin{document}
\section*{UnsolvedMath ID: 30001823}
\textbf{Source problem code:} catalogue entry
\subsection{Definitions and mathematical statement}
A metric space $(X,d)$ is geodesic if, for every $p,q\in X$, there is an isometric map $\gamma:[0,d(p,q)]\to X$ with $\gamma(0)=p$ and $\gamma(d(p,q))=q$. Fix a basepoint $x_0\in X$. A pointed isometry is a bijective isometry that sends the distinguished basepoint to the distinguished basepoint.

The supplied question concerns tangents obtained by blowing up the metric near $x_0$, but it does not specify a notion of convergence. One customary interpretation, stated here conditionally, uses pointed Gromov--Hausdorff convergence of proper metric spaces. A metric space is proper if every closed bounded ball is compact. If $X$ is additionally assumed proper, define
\[
\operatorname{Tan}_{\mathrm{pGH}}(X,x_0)
=
\left\{
[(Y,d_Y,y_0)]_{\mathrm{pt}}:
\begin{array}{l}
(Y,d_Y)\text{ is a proper metric space, and there exists}\\
(r_j)_{j\in\N}\text{ with }r_j>0,\ r_j\to0,\text{ such that}\\
(X,r_j^{-1}d,x_0)\xrightarrow{\mathrm{pGH}}(Y,d_Y,y_0)
\end{array}
\right\},
\]
where $[\cdot]_{\mathrm{pt}}$ denotes the pointed-isometry class. Under this interpretation, uniqueness means
\[
\bigl|\operatorname{Tan}_{\mathrm{pGH}}(X,x_0)\bigr|=1,
\]
so it includes existence of a tangent. No existence or subsequential compactness assumption is supplied in the record.

\textbf{Source question.} For a geodesic metric space $(X,d)$ and $x_0\in X$, how large can the collection $\operatorname{Tan}(X,x_0)$ of pointed-isometry classes of tangents be, and what structural consequences follow when this collection has exactly one element?

\textbf{Unresolved qualifications.} The proper-space pointed Gromov--Hausdorff definition above is a possible interpretation, not a hypothesis or definition established by the supplied source record. The original question assumes only geodesicity. It does not specify convergence for nonproper spaces, a cardinal bound or classification to establish, or the structural conclusions intended by uniqueness. Consequently, the supplied question does not determine a single precise conjecture without further source clarification.
\subsection{Short English statement}
For a geodesic metric space with a chosen basepoint, ask how many distinct tangents can arise and what the existence of exactly one tangent implies. A conditional pointed Gromov--Hausdorff formulation is given, but the source leaves the convergence convention, additional hypotheses, and intended rigidity conclusions unspecified.
\subsection{Sources}
\begin{itemize}
\item Supplied local UnsolvedMath registry record, permanent ID \texttt{30001823}, title \emph{Cardinality and Rigidity of Metric Tangent Spaces}, source problem code \texttt{catalogue entry}, in the collection \emph{Oberwolfach Reports Open Problems}.
\item Supplied report citation: \emph{Differentialgeometrie im Großen (2012)}, Oberwolfach Reports, pp.~1857--1912, DOI \texttt{10.4171/owr/2011/33}. Source URL and supplied literature reference labelled ``original Oberwolfach report'': \url{https://doi.org/10.4171/owr/2011/33}. The background describes these reports as records of mathematical workshops and their open-problem discussions. No individual proposer is supplied.
\item The registry gives the proposed year as 2011, while the supplied report title contains 2012. These details are retained without resolving their relationship.
\item Literature assessment dated exactly \texttt{2026-08-21}: the supplied curation labels the problem open and describes it as broad questions about tangent cones and uniqueness of blow-ups, with no resolution supplied and openness asserted at this level of generality. This is supplied dated curation derived from the OWR record, not an independently verified literature assessment; time-sensitive status and cited claims require independent verification.
\item The supplied statement status is \texttt{corrected\_verified}; its accompanying note reports manual semantic, typographic, and LaTeX review of a meaning-preserving cleanup. This records the supplied editorial assessment and does not establish mathematical correctness or resolve the formulation ambiguities.
\end{itemize}
\end{document}
